% SGA 3, Expose I, Loop 1 English reconstruction.
% Scope: Section 1.6 through the end of Remark 1.7.3.
% Authority: Polo--Gille Exp1-13oct24.pdf, component pages 5--10
% (combined-reader pages 15--20). OCR is drafting/locator material only.

\subsection*{\SGAInlineRefTarget{sga3:i:subsection:1-6}1.6. 函子中的基变换}

对每个 \(S\in\operatorname{Ob}(\mathscr C)\), 有典范函子
\[
  i_S:\mathscr C_{/S}\longrightarrow\mathscr C
\]
其定义为: 当 \(f\) 是箭头 \(T\to S\) 时, \(i_S(f)=T\). 若
\(f:T\to S\) 是 \(\mathscr C_{/S}\) 的一个对象, 我们用
\(i_{T/S}=i_f\) 表示函子
\[
  i_{T/S}:(\mathscr C_{/S})_{/T}\longrightarrow\mathscr C_{/S},
\]
并且有交换图:

% Authority: Exp1-13oct24.pdf p. 5; combined reader p. 15.
\begin{center}
  \begin{tikzcd}[column sep=4.7em,row sep=3.2em]
  (\mathscr C_{/S})_{/T}
    \arrow[r,"i_{T/S}"]
    \arrow[dr,"\simeq"']
  &
  \mathscr C_{/S}
    \arrow[r,"i_S"]
  &
  \mathscr C
  \\
  &
  \mathscr C_{/T}
    \arrow[ur,"i_T"']
  &
  {} ,
\end{tikzcd}

\end{center}

因此, 将 \((\mathscr C_{/S})_{/T}\) 与 \(\mathscr C_{/T}\) 等同
(今后我们都这样做), 就有
\[
  i_S\circ i_{T/S}=i_T.
\]
同样, 当 \(\mathscr C\) 有终对象 \(e\), 且将 \(\mathscr C\) 与
\(\mathscr C_{/e}\) 等同时, \(i_{S/e}:\mathscr C_{/S}\to\mathscr C_{/e}\)
便与 \(i_S\) 等同.

对 \(X\in\operatorname{Ob}(\mathscr C)\) (分别地,
\(Y\in\operatorname{Ob}(\mathscr C_{/S})\)), 用 \(p_S(X)\)
(分别地, \(p_{T/S}(Y)\)) 表示 \(\mathscr C_{/S}\) (分别地,
\(\mathscr C_{/T}\)) 中如下定义的对象 (只要它存在):
\(X\times S\) (分别地, \(Y\times_S T\)) 配备其第二投影:

% Authority: Exp1-13oct24.pdf p. 5; combined reader p. 15.
\begin{center}
  \begin{tikzcd}[column sep=6.5em,row sep=3.2em]
  X\times S
    \arrow[d,"p_S(X)"']
  &
  \textup{resp.}
  &
  Y\times_S T
    \arrow[d,"p_{T/S}(Y)"]
  \\
  S
  &&
  T
  \qquad .
\end{tikzcd}

\end{center}

部分定义的函子 \(p_S\) (分别地, \(p_{T/S}\)) 称为
\emph{基变换函子 (base-change functor)}. 由积 (分别地, 纤维积) 的定义,
它是 \(i_S\) (分别地, \(i_{T/S}\)) 的右伴随.
\footnote{编者注: 即若 \(g:U\to S\) (分别地,
\(h:V\to T\)) 是 \(\mathscr C_{/S}\) (分别地,
\(\mathscr C_{/T}\)) 的一个对象, 则 \(i_S(g)=U\), 而 \(i_{T/S}(h)\)
是 \(\mathscr C_{/S}\) 中的对象 \(f\circ h:V\to S\), 并且
\[
  \operatorname{Hom}_{\mathscr C_{/S}}(U,X\times S)
    \simeq \operatorname{Hom}_{\mathscr C}(U,X),
\]
分别地,
\[
  \operatorname{Hom}_{\mathscr C_{/T}}(V,X\times_S T)
    \simeq \operatorname{Hom}_{\mathscr C_{/S}}(V,X).
\]}
我们也记
\[
  p_S(X)=X_S,
  \qquad
  p_{T/S}(Y)=Y_T.
\]

函子 \(i_S\) 定义了一个限制函子
\[
  i_S^*:\widehat{\mathscr C}\longrightarrow
    \widehat{\mathscr C_{/S}}.
\]
我们记
\[
  \mathbf F_S=i_S^*(\mathbf F)=\mathbf F\circ i_S.
\]
显然,
\[
  i_{T/S}^*\circ i_S^*=i_T^*,
\]
即对每个函子
\(\mathbf F\in\operatorname{Ob}(\widehat{\mathscr C})\),
\[
  (\mathbf F_S)_T=\mathbf F_T.
\]
下面的结果说明这一记号是合理的.

\medskip
\noindent\SGAInlineRefTarget{sga3:i:proposition:1-6-1}\textbf{命题 1.6.1.}
\textit{函子
\[
  (\mathbf h_X)_S:(\mathscr C_{/S})^{\circ}\longrightarrow\mathbf{Set}
\]
可表当且仅当积 \(X\times S\) 存在. 此时
\[
  (\mathbf h_X)_S\simeq\mathbf h_{X_S}.
\]}

这说明 \(\mathbf F_S\) 有两种解释: 它既是 \(\mathbf F\) 在
\(\mathscr C_{/S}\) 上的限制, 也是由基变换 \(e\leftarrow S\)
得到的函子. 因此采用如下记号:

% Authority: Exp1-13oct24.pdf p. 6; combined reader p. 16.
\begin{center}
  \begin{tikzcd}[column sep=3.8em,row sep=3.0em]
  \mathbf F
    \arrow[r,no head]
    \arrow[d,no head]
  &
  \mathbf F_S
    \arrow[r,no head]
    \arrow[d,no head]
  &
  \mathbf F_T
    \arrow[d,no head]
  \\
  e
  &
  S
    \arrow[l]
  &
  T
    \arrow[l]
\end{tikzcd}

\end{center}

此图同时体现了前面的两种解释. 注意
\[
  \Gamma(\mathbf F_S)
    \simeq\operatorname{Hom}(\mathbf h_S,\mathbf F)
    \simeq\mathbf F(S),
\]
特别地,
\[
  \Gamma(X_S)\simeq\operatorname{Hom}(S,X).
\]

\subsection*{\SGAInlineRefTarget{sga3:i:subsection:1-7-0}1.7.0.}

设 \(\mathbf E\) 是 \(\widehat{\mathscr C}\) 的一个对象.
\footnote{编者注: 下面的引理是新增的 (见 SGA~4,
I.3.4); 它用于 \SGARefLink{sga3:i:proposition:1-7-1}{1.7.1} 的证明,
并且在后文中会多次用到.}
考虑 \(\mathbf E\) 上的 \emph{\(\mathscr C\) 的对象}所成的范畴
\(\mathscr C_{/\mathbf E}\). 其对象是二元组 \((V,\rho)\),
由 \(\mathscr C\) 的一个对象 \(V\) 和 \(\widehat{\mathscr C}\) 中的一个态射
\(\rho:\mathbf h_V\to\mathbf E\), 即 \(\rho\in\mathbf E(V)\) 组成.
从 \((V,\rho)\) 到 \((V',\rho')\) 的态射是 \(\mathscr C\) 中满足
\[
  \rho'\circ f=\rho,
  \qquad\text{即}\qquad
  \mathbf E(f)(\rho')=\rho.
\]
的态射 \(f:V\to V'\).
设 \(\mathbf L\) 为函子
\[
  \mathbf L=\varinjlim_{(V,\rho)\in\mathscr C_{/\mathbf E}}\mathbf h_V.
\]
因此, 对每个 \(S\in\operatorname{Ob}(\mathscr C)\),
\[
  \mathbf L(S)
    =\varinjlim_{(V,\rho)\in\mathscr C_{/\mathbf E}}\mathbf h_V(S)
\]
是三元组 \((V,\rho,v)\) 的等价类所成的集合, 其中
\(v:S\to V\) 是 \(\mathscr C\) 中的态射; 对每个满足
\(\rho'\circ f=\rho\) 的 \(\mathscr C\) 中的态射 \(f:V\to V'\),
将 \((V,\rho,v)\) 与 \((V',\rho',f\circ v)\) 等同.

将 \((V,\rho,v)\) 的类映到 \(\mathbf E(S)\) 中的元素
\(\rho\circ v\) 的映射 \(\phi_{\mathbf E}(S)\) 是良定义的,
并给出一个函子态射
\[
  \phi_{\mathbf E}:
  \varinjlim_{(V,\rho)\in\mathscr C_{/\mathbf E}}\mathbf h_V
    \longrightarrow\mathbf E.
\]

\medskip
\noindent\textbf{引理.}
\textit{态射 \(\phi_{\mathbf E}\) 是同构.}

事实上, 设 \(S\in\operatorname{Ob}(\mathscr C)\). 每个
\(x\in\mathbf E(S)\) 都是三元组 \((S,x,\operatorname{id}_S)\)
在 \(\phi_{\mathbf E}(S)\) 下的像; 因此 \(\phi_{\mathbf E}(S)\) 是满射.
反过来, 设
\(\ell_1=(V_1,\rho_1,v_1)\) 和
\(\ell_2=(V_2,\rho_2,v_2)\) 是 \(\mathbf L(S)\) 的两个元素,
它们在 \(\mathbf E(S)\) 中具有相同的像, 并置
\[
  \gamma=\rho_1\circ v_1=\rho_2\circ v_2.
\]
那么 \(\ell_1\) 和 \(\ell_2\) 在 \(\mathbf L(S)\) 中都等于
\((S,\gamma,\operatorname{id}_S)\) 的类. 因此
\(\phi_{\mathbf E}(S)\) 是单射.

\medskip
\noindent\textbf{推论.}
\textit{对 \(\widehat{\mathscr C}\) 的每个对象 \(\mathbf F\), 有
\[
  \operatorname{Hom}(\mathbf E,\mathbf F)
    =\varprojlim_{(V,\rho)\in\mathscr C_{/\mathbf E}}\mathbf F(V).
\]}

\SGARefTarget{sga3:i:subsection:1:7}
\subsection*{1.7. 对象 \(\underline{\operatorname{Hom}}\),
\(\underline{\operatorname{Isom}}\) 等}

设 \(\mathbf F\) 和 \(\mathbf G\) 是 \(\widehat{\mathscr C}\)
的两个对象. 定义 \(\widehat{\mathscr C}\) 的另一个对象如下:
\begin{align*}
  \underline{\operatorname{Hom}}(\mathbf F,\mathbf G)(S)
  &={}\operatorname{Hom}_{\widehat{\mathscr C_{/S}}}
       (\mathbf F_S,\mathbf G_S)\\
  &\simeq\operatorname{Hom}_{\widehat{\mathscr C}_{/\mathbf h_S}}
       (\mathbf F\times\mathbf h_S,\mathbf G\times\mathbf h_S)\\
  &\simeq\operatorname{Hom}_{\widehat{\mathscr C}}
       (\mathbf F\times\mathbf h_S,\mathbf G).
\end{align*}
对象 \(\underline{\operatorname{Hom}}(\mathbf F,\mathbf G)\) 有如下性质:
\begin{enumerate}
  \item[(i)]
    \(\underline{\operatorname{Hom}}(\underline{\mathbf e},\mathbf G)
      \simeq\mathbf G\).
    \footnote{编者注: 若 \(\mathbf E\) 是 \(\widehat{\mathscr C}\)
    的第三个对象, 则
    \[
      \underline{\operatorname{Hom}}(\mathbf E,\mathbf F\times\mathbf G)
      \simeq
      \underline{\operatorname{Hom}}(\mathbf E,\mathbf F)
        \times
      \underline{\operatorname{Hom}}(\mathbf E,\mathbf G).
    \]}
  \item[\SGAInlineRefTarget{sga3:i:subsection:1:7:item:ii}(ii)] \(\underline{\operatorname{Hom}}\)
    的构造与基扩张可交换:
    \[
      \underline{\operatorname{Hom}}(\mathbf F_S,\mathbf G_S)
        \simeq
      \underline{\operatorname{Hom}}(\mathbf F,\mathbf G)_S.
    \]
  \item[(iii)]
    \((\mathbf F,\mathbf G)\mapsto
      \underline{\operatorname{Hom}}(\mathbf F,\mathbf G)\) 是双函子,
    对 \(\mathbf F\) 反变, 对 \(\mathbf G\) 协变.
\end{enumerate}
这三个性质由定义立即可得.

我们将证明, 对 \(\widehat{\mathscr C}\) 的每个对象 \(\mathbf E\),
\[
  \operatorname{Hom}
    \bigl(\mathbf E,
      \underline{\operatorname{Hom}}(\mathbf F,\mathbf G)\bigr)
    \simeq
  \operatorname{Hom}(\mathbf E\times\mathbf F,\mathbf G).
\]
设 \(\phi:\mathbf E\times\mathbf F\to\mathbf G\). 我们要为它配一个
从 \(\mathbf E\) 到 \(\underline{\operatorname{Hom}}(\mathbf F,\mathbf G)\)
的态射. 对 \(\mathscr C\) 的一个箭头 \(S'\to S\), 有映射
\[
  \mathbf E(S)\times\mathbf F(S')
    \longrightarrow
  \mathbf E(S')\times\mathbf F(S')
    \xrightarrow{\ \phi(S')\ }
  \mathbf G(S').
\]
因此, \(\mathbf E(S)\) 的每个元素 \(e\) 对每个 \(S'\to S\)
定义一个关于 \(S'\) 函子性的映射 \(\mathbf F(S')\to\mathbf G(S')\);
这便是 \(\underline{\operatorname{Hom}}(\mathbf F,\mathbf G)(S)\)
的一个元素 \(\theta_\phi(e)\). 由此得到映射
\[
  \phi\longmapsto\theta_\phi,
  \qquad
  \operatorname{Hom}(\mathbf E\times\mathbf F,\mathbf G)
    \longrightarrow
  \operatorname{Hom}
    \bigl(\mathbf E,
      \underline{\operatorname{Hom}}(\mathbf F,\mathbf G)\bigr),
\]
它关于 \(\mathbf E\) 是函子性的.

\medskip
\noindent\SGAInlineRefTarget{sga3:i:proposition:1-7-1}\textbf{命题 1.7.1.}
\footnote{编者注: 新增的 \SGARefLink{sga3:i:proposition:1-7-1:item:b}{(b)}
项将在 II.1 和 II.3.11 中用到. 还补写了
\SGARefLink{sga3:i:proposition:1-7-1:item:a}{(a)} 项的另一种更直接的证明梗概.}
\textit{设
\(\mathbf E,\mathbf F,\mathbf G\in
\operatorname{Ob}(\widehat{\mathscr C})\).
\begin{enumerate}
  \item[\SGAInlineRefTarget{sga3:i:proposition:1-7-1:item:a}(a)] 映射 \(\phi\mapsto\theta_\phi\) 是双射:
    \[
      \operatorname{Hom}(\mathbf E\times\mathbf F,\mathbf G)
        \xrightarrow{\ \sim\ }
      \operatorname{Hom}
        \bigl(\mathbf E,
          \underline{\operatorname{Hom}}(\mathbf F,\mathbf G)\bigr).
    \]
  \item[\SGAInlineRefTarget{sga3:i:proposition:1-7-1:item:b}(b)] 此外, 有函子同构
    \[
      \underline{\operatorname{Hom}}
        \bigl(\mathbf E,
          \underline{\operatorname{Hom}}(\mathbf F,\mathbf G)\bigr)
        \simeq
      \underline{\operatorname{Hom}}
        (\mathbf E\times\mathbf F,\mathbf G).
    \]
\end{enumerate}}

对于 \SGARefLink{sga3:i:proposition:1-7-1:item:a}{(a)}, 将两端都看成
\(\mathbf E\) 的函子. 当 \(\mathbf E=\mathbf h_X\) 时, 所述结果成立,
因为这时它恰好是函子
\(\underline{\operatorname{Hom}}(\mathbf F,\mathbf G)\) 的定义. 此外,
两端作为 \(\mathbf E\) 的函子, 都将正向极限变成逆向极限.
最后, 由 \SGARefLink{sga3:i:subsection:1-7-0}{引理~1.7.0},
\(\widehat{\mathscr C}\) 的每个对象 \(\mathbf E\) 都同构于
\(\mathbf h_X\) 的正向极限, 其中 \(X\) 遍历 \(\mathscr C_{/\mathbf E}\).
这就证明了 \SGARefLink{sga3:i:proposition:1-7-1:item:a}{(a)}.

我们还给出 \SGARefLink{sga3:i:proposition:1-7-1:item:a}{(a)} 的一个直接证明的梗概.
对每个
\[
  \theta\in\operatorname{Hom}
    \bigl(\mathbf E,
      \underline{\operatorname{Hom}}(\mathbf F,\mathbf G)\bigr)
\]
如下配一个 \(\phi_\theta\in
\operatorname{Hom}(\mathbf E\times\mathbf F,\mathbf G)\). 对每个
\(S\in\operatorname{Ob}(\mathscr C)\), 有关于 \(S\) 函子性的映射
\[
  \theta(S):\mathbf E(S)\longrightarrow
    \operatorname{Hom}(\mathbf F\times S,\mathbf G),
\]
若 \((e,f)\in\mathbf E(S)\times\mathbf F(S)\), 则 \(f\)
是态射 \(S\to\mathbf F\), 所以 \(f\times\operatorname{id}_S\)
是态射 \(S\to\mathbf F\times S\). 另一方面, \(\theta(S)(e)\)
是态射 \(\mathbf F\times S\to\mathbf G\). 因此复合给出态射
\[
  \theta(S)(e)\circ(f\times\operatorname{id}_S):S\longrightarrow\mathbf G,
\]
即 \(\mathbf G(S)\) 的一个元素 \(\phi_\theta(S)(e,f)\). 容易验证,
\(S\mapsto\phi_\theta(S)\) 关于 \(S\) 是函子性的, 因而定义了
从 \(\mathbf E\times\mathbf F\) 到 \(\mathbf G\) 的态射 \(\phi_\theta\).
映射 \(\theta\mapsto\phi_\theta\) 和 \(\phi\mapsto\theta_\phi\)
是互逆双射的验证留给读者.

对于 \SGARefLink{sga3:i:proposition:1-7-1:item:b}{(b)}, 若
\(S\in\operatorname{Ob}(\mathscr C)\), 则由
\SGARefLink{sga3:i:subsection:1:7:item:ii}{1.7(ii)}, 以及将
\SGARefLink{sga3:i:proposition:1-7-1:item:a}{(a)} 项应用于 \(\mathscr C_{/S}\), 有
\begin{align*}
  &\underline{\operatorname{Hom}}
    \bigl(\mathbf E,
      \underline{\operatorname{Hom}}(\mathbf F,\mathbf G)\bigr)(S)\\
  &\quad\simeq
    \operatorname{Hom}_S
      \bigl(\mathbf E_S,
        \underline{\operatorname{Hom}}_S(\mathbf F_S,\mathbf G_S)\bigr)\\
  &\quad\simeq
    \operatorname{Hom}_S
      (\mathbf E_S\times_S\mathbf F_S,\mathbf G_S)\\
  &\quad\simeq
    \operatorname{Hom}(\mathbf E\times\mathbf F\times S,\mathbf G)\\
  &\quad\simeq
    \underline{\operatorname{Hom}}
      (\mathbf E\times\mathbf F,\mathbf G)(S),
\end{align*}
并且这些同构关于 \(S\) 是函子性的.

\medskip
\noindent\SGAInlineRefTarget{sga3:i:corollary:1-7-2}\textbf{推论 1.7.2.}
\textit{有同构
\[
  \operatorname{Hom}
    \bigl(\mathbf E,
      \underline{\operatorname{Hom}}(\mathbf F,\mathbf G)\bigr)
    \simeq
  \operatorname{Hom}
    \bigl(\mathbf F,
      \underline{\operatorname{Hom}}(\mathbf E,\mathbf G)\bigr),
\]
以及
\[
  \underline{\operatorname{Hom}}
    \bigl(\mathbf E,
      \underline{\operatorname{Hom}}(\mathbf F,\mathbf G)\bigr)
    \simeq
  \underline{\operatorname{Hom}}
    \bigl(\mathbf F,
      \underline{\operatorname{Hom}}(\mathbf E,\mathbf G)\bigr).
\]}

特别地, 取 \(\mathbf E=\underline{\mathbf e}\), 并利用
\(\underline{\operatorname{Hom}}(\underline{\mathbf e},\mathbf G)
\simeq\mathbf G\), 得到
\[
  \Gamma\bigl(\underline{\operatorname{Hom}}(\mathbf F,\mathbf G)\bigr)
    \simeq\operatorname{Hom}(\mathbf F,\mathbf G).
\]
复合运算给出函子性的态射
\[
  \underline{\operatorname{Hom}}(\mathbf F,\mathbf G)
    \times
  \underline{\operatorname{Hom}}(\mathbf G,\mathbf H)
    \longrightarrow
  \underline{\operatorname{Hom}}(\mathbf F,\mathbf H).
\]

若 \(\mathbf F\) 和 \(\mathbf G\) 是 \(\widehat{\mathscr C}\)
的两个对象, 用 \(\operatorname{Isom}(\mathbf F,\mathbf G)\) 表示
\(\operatorname{Hom}(\mathbf F,\mathbf G)\) 中由从 \(\mathbf F\)
到 \(\mathbf G\) 的同构组成的子集. 定义
\(\underline{\operatorname{Hom}}(\mathbf F,\mathbf G)\) 的子对象
\(\underline{\operatorname{Isom}}(\mathbf F,\mathbf G)\) 如下:
\[
  \underline{\operatorname{Isom}}(\mathbf F,\mathbf G)(S)
    =\operatorname{Isom}(\mathbf F_S,\mathbf G_S).
\]
则
\[
  \Gamma\bigl(\underline{\operatorname{Isom}}
    (\mathbf F,\mathbf G)\bigr)
    \simeq\operatorname{Isom}(\mathbf F,\mathbf G),
  \qquad
  \underline{\operatorname{Isom}}(\mathbf F,\mathbf G)
    \simeq
  \underline{\operatorname{Isom}}(\mathbf G,\mathbf F).
\]
在 \(\mathbf F=\mathbf G\) 的特殊情形下, 置
\begin{align*}
  \underline{\operatorname{End}}(\mathbf F)
    &=\underline{\operatorname{Hom}}(\mathbf F,\mathbf F),
  &
  \operatorname{End}(\mathbf F)
    &=\operatorname{Hom}(\mathbf F,\mathbf F)
      \simeq\Gamma\bigl(\underline{\operatorname{End}}(\mathbf F)\bigr),\\
  \underline{\operatorname{Aut}}(\mathbf F)
    &=\underline{\operatorname{Isom}}(\mathbf F,\mathbf F),
  &
  \operatorname{Aut}(\mathbf F)
    &=\operatorname{Isom}(\mathbf F,\mathbf F)
      \simeq\Gamma\bigl(\underline{\operatorname{Aut}}(\mathbf F)\bigr).
\end{align*}
对象 \(\underline{\operatorname{Hom}}\),
\(\underline{\operatorname{Isom}}\),
\(\underline{\operatorname{Aut}}\) 和
\(\underline{\operatorname{End}}\) 的构造都与基变换可交换.

\medskip
\noindent\SGAInlineRefTarget{sga3:i:remark:1-7-3}\textbf{注 1.7.3.}
\footnote{编者注: 为便于后文引用, 增加了编号
\SGARefLink{sga3:i:remark:1-7-3}{1.7.3}.}
也可以如下构造一个与
\(\underline{\operatorname{Isom}}(\mathbf F,\mathbf G)\) 同构的对象.
有态射
\[
  \underline{\operatorname{Hom}}(\mathbf F,\mathbf G)
    \times
  \underline{\operatorname{Hom}}(\mathbf G,\mathbf F)
    \longrightarrow
  \underline{\operatorname{End}}(\mathbf F).
\]
交换 \(\mathbf F\) 和 \(\mathbf G\), 得到态射
\[
  \underline{\operatorname{Hom}}(\mathbf F,\mathbf G)
    \times
  \underline{\operatorname{Hom}}(\mathbf G,\mathbf F)
    \longrightarrow
  \underline{\operatorname{End}}(\mathbf F)
    \times
  \underline{\operatorname{End}}(\mathbf G).
\]
另一方面, \(\mathbf F\) 的恒等态射是
\(\operatorname{End}(\mathbf F)\) 的一个元素, 因而定义了态射
\(\underline{\mathbf e}\to
\underline{\operatorname{End}}(\mathbf F)\). 对 \(\mathbf G\) 作同样处理,
再取积, 得到态射
\[
  \underline{\mathbf e}\longrightarrow
  \underline{\operatorname{End}}(\mathbf F)
    \times
  \underline{\operatorname{End}}(\mathbf G).
\]
于是立即可见, \(\underline{\mathbf e}\) 与
\[
  \underline{\operatorname{Hom}}(\mathbf F,\mathbf G)
    \times
  \underline{\operatorname{Hom}}(\mathbf G,\mathbf F)
\]
在
\(\underline{\operatorname{End}}(\mathbf F)
\times\underline{\operatorname{End}}(\mathbf G)\)
上的纤维积同构于
\(\underline{\operatorname{Isom}}(\mathbf F,\mathbf G)\).

所有这些定义特别适用于
\(\mathbf F=\mathbf h_X\) 和 \(\mathbf G=\mathbf h_Y\) 的情形.
若 \(\underline{\operatorname{Hom}}(\mathbf h_X,\mathbf h_Y)\)
可由 \(\mathscr C\) 的某个对象表示, 则将该对象记为
\(\underline{\operatorname{Hom}}(X,Y)\). 它具有如下性质:
若 \(Z\times X\) 存在, 则
\[
  \operatorname{Hom}
    \bigl(Z,\underline{\operatorname{Hom}}(X,Y)\bigr)
    \simeq
  \operatorname{Hom}(Z\times X,Y).
\]
当 \(\mathscr C\) 中的积存在时, 这一性质刻画了该对象.

同样, 当下列对象存在时, 定义
\[
  \underline{\operatorname{Isom}}(X,Y),
  \qquad
  \underline{\operatorname{End}}(X),
  \qquad
  \underline{\operatorname{Aut}}(X).
\]
我们只指出, 上述构造表明, 只要 \(\mathscr C\) 中的纤维积存在,
并且对象
\[
  \underline{\operatorname{Hom}}(X,Y),
  \quad
  \underline{\operatorname{Hom}}(Y,X),
  \quad
  \underline{\operatorname{End}}(X),
  \quad\text{以及}\quad
  \underline{\operatorname{End}}(Y)
\]
存在, 则 \(\underline{\operatorname{Isom}}(X,Y)\) 存在.

以上一切也适用于形如 \(\mathscr C_{/S}\) 的范畴.
我们将尽可能用适当的符号明确表示相应的对象. 例如, 若 \(T\) 和
\(T'\) 是 \(\mathscr C\) 中 \(S\) 上的两个对象, 则用
\(\underline{\operatorname{Hom}}_S(T,T')\) 表示对象
\[
  \underline{\operatorname{Hom}}_{\mathscr C_{/S}}(T/S,T'/S).
\]

% End of assigned component. Section 1.8 begins in the next unit.

% Translation coverage: entire component, Sections 1.6--1.7.3.
% Source lines: 533; translated body lines before this audit comment: 518.
% Math: 191 inline, 44 displayed, 3 align* blocks; all preserved (two prose \text arguments translated).
% Footnotes: 5/5; reference targets: 10/10; reference links: 11/11; unresolved TODOs: 0.
